\documentclass[../main.tex]{subfiles}
\usepackage{setspace,longtable}
\setstretch{1.5}
\begin{document}

\section*{DIAGNÓSTICO AULICO: AÑO 2026}
\begin{tabular}{r l}
	\textbf{CURSO :} & 4° 1° \\
\end{tabular}

\subsection*{Descripción de la propuesta de diagnóstico}
% [Insert description]

\subsection*{Instrumentos utilizados}
% [Insert instruments]

\subsection*{DESCRIPCIÓN DE RESULTADOS OBTENIDOS}

\subsubsection*{Información cuantitativa}
\begin{longtable}{|p{4cm}|p{3cm}|p{3cm}|p{3cm}|}
	\hline
	Análisis y comprensión & No logrado & Logrado & Altamente logrado \\ \hline
	                       &            &         &                   \\ \hline
\end{longtable}

\subsubsection*{Información cualitativa}
\begin{itemize}
	\item El grupo en relación consigo mismo:
	\item El grupo ante los hábitos:
	\item El grupo ante los aprendizajes:
	\item El grupo en relación con el material de estudio:
	\item El grupo en relación con el docente:
	\item El grupo y las necesidades particulares:
\end{itemize}

\subsection*{Método evaluativo sugerido}
% [Insert method]

\subsection*{Estudiantes que presentan dificultades}
\begin{longtable}{|p{6cm}|p{6cm}|}
	\hline
	Alumnos & Dificultad \\ \hline
	        &            \\ \hline
\end{longtable}

\end{document}
